% GENERATED FILE. Edit canonical lessons, then run npm run generate:books.
% canonical-source-hash: fnv1a64:f81e3f9b60c66f6a
% canonical-lessons: ZH-C415-hangzhanlou, ZH-C415-tuoyun, ZH-C415-dengjipai, ZH-C415-anjian, ZH-C415-dengjikou

\chapter{The Terminal}
\label{ch:zh-the-terminal}
\hlchaptermodality{415}

\begin{chapteropening}
\textbf{By the end of this chapter:} \emph{I can find the terminal, check in a bag, show my boarding pass, go through security and find the gate.}

Complete the last lesson of chapter 415: I can find the terminal, check in a bag, show my boarding pass, go through security and find the gate.
\end{chapteropening}

\section[hángzhànlóu]{hángzhànlóu — airport terminal}

\label{lesson:ZH-C415-hangzhanlou}



\begin{quote}
Before the new one: say the Mandarin for to get airsick, then the Mandarin for a seat belt.
\end{quote}

\subsection*{You'll want to know: hángzhànlóu}
\emph{hángzhànlóu} — \textquotedblleft{}airport terminal\textquotedblright{}.

\begin{quote}
\emph{wǒmen de hángbān zài èr hào hángzhànlóu} — Our flight goes from terminal two.
\end{quote}

Its characters come later; for now it is a word for the ear.

\begin{cousinweb}
\emph{hángzhànlóu} is built from \emph{háng} 'to sail, to fly', \zh{站} 'a station' and \emph{lóu} 'a building', all inherited from Middle Chinese, \emph{hang}, \emph{treamH} and \emph{luw} (Guangyun, in Baxter's transcription), which glosses the first 'a boat' and the third as a building of several storeys.
\end{cousinweb}

\begin{grammarlens}[title={What you've built}]
One more word for getting around.
\end{grammarlens}

\subsection*{Your turn}
\begin{itemize}
\raggedright
  \item \emph{Say it:} \emph{hángzhànlóu}
  \item \emph{Say it:} \emph{hángzhànlóu}, once more
  \item \emph{Say it:} \emph{wǒmen de hángbān zài èr hào hángzhànlóu} — Our flight goes from terminal two.
  \item \emph{Recall:} say the Mandarin for to get airsick, then the Mandarin for a seat belt, then say \emph{hángzhànlóu} again
\end{itemize}

\begin{tcolorbox}[breakable,colback=teal!4,colframe=teal!35!black,title={Before you move on}]
What does hángzhànlóu mean? (\textquotedblleft{}Airport terminal\textquotedblright{}.) Say it once more.
\end{tcolorbox}
\section[tuōyùn]{tuōyùn — to check in a suitcase; to send luggage in the hold}

\label{lesson:ZH-C415-tuoyun}



\begin{quote}
Before the new one: say the Mandarin for a seat belt, then the Mandarin for an airport terminal building.
\end{quote}

\subsection*{You'll want to know: tuōyùn}
\emph{tuōyùn} — \textquotedblleft{}to check in a suitcase; to send luggage in the hold\textquotedblright{}.

\begin{quote}
\emph{zhège xiāngzi yào tuōyùn} — This suitcase needs to be checked in.
\end{quote}

Its characters come later; for now it is a word for the ear.

\begin{cousinweb}
\emph{tuōyùn} is built from \emph{tuō} 'to entrust' and \emph{yùn} 'to carry', both inherited from Middle Chinese, \emph{thak} and \emph{hjunH} (Guangyun, in Baxter's transcription), which lists the first under another spelling, 'to entrust', and glosses the second 'to move, to carry'.
\end{cousinweb}

\begin{grammarlens}[title={What you've built}]
One more word for getting around.
\end{grammarlens}

\subsection*{Your turn}
\begin{itemize}
\raggedright
  \item \emph{Say it:} \emph{tuōyùn}
  \item \emph{Say it:} \emph{tuōyùn}, once more
  \item \emph{Say it:} \emph{zhège xiāngzi yào tuōyùn} — This suitcase needs to be checked in.
  \item \emph{Recall:} say the Mandarin for a seat belt, then the Mandarin for an airport terminal building, then say \emph{tuōyùn} again
\end{itemize}

\begin{tcolorbox}[breakable,colback=teal!4,colframe=teal!35!black,title={Before you move on}]
What does tuōyùn mean? (\textquotedblleft{}To check in a suitcase; to send luggage in the hold\textquotedblright{}.) Say it once more.
\end{tcolorbox}
\section[dēngjīpái]{dēngjīpái — boarding pass}

\label{lesson:ZH-C415-dengjipai}



\begin{quote}
Before the new one: say the Mandarin for an airport terminal building, then the Mandarin for to check in a suitcase.
\end{quote}

\subsection*{You'll want to know: dēngjīpái}
\emph{dēngjīpái} — \textquotedblleft{}boarding pass\textquotedblright{}.

\begin{quote}
\emph{wǒ de dēngjīpái bú jiàn le} — My boarding pass has gone missing.
\end{quote}

Its characters come later; for now it is a word for the ear.

\begin{cousinweb}
\emph{dēngjīpái} is built from \emph{dēngjī} 'to board a plane' and \emph{pái} 'a card', all inherited from Middle Chinese, \emph{tong}, \emph{kj+j} and \emph{bea} (Guangyun, in Baxter's transcription), which glosses the last as a signboard.
\end{cousinweb}

\begin{grammarlens}[title={What you've built}]
One more word for getting around.
\end{grammarlens}

\subsection*{Your turn}
\begin{itemize}
\raggedright
  \item \emph{Say it:} \emph{dēngjīpái}
  \item \emph{Say it:} \emph{dēngjīpái}, once more
  \item \emph{Say it:} \emph{wǒ de dēngjīpái bú jiàn le} — My boarding pass has gone missing.
  \item \emph{Recall:} say the Mandarin for an airport terminal building, then the Mandarin for to check in a suitcase, then say \emph{dēngjīpái} again
\end{itemize}

\begin{tcolorbox}[breakable,colback=teal!4,colframe=teal!35!black,title={Before you move on}]
What does dēngjīpái mean? (\textquotedblleft{}Boarding pass\textquotedblright{}.) Say it once more.
\end{tcolorbox}
\section[\texorpdfstring{\zh{安检} (ānjiǎn)}{安检 (ānjiǎn)}]{\zh{安检} (ānjiǎn) — security check}

\label{lesson:ZH-C415-anjian}



\begin{quote}
Before the new one: say the Mandarin for to check in a suitcase, then the Mandarin for a boarding pass.
\end{quote}

\subsection*{You'll want to know: \zh{安检}}
\textbf{\zh{安检}} — \emph{ānjiǎn} — \textquotedblleft{}security check\textquotedblright{}.

\begin{quote}
\textbf{\zh{我们先去安检}} \emph{wǒmen xiān qù ānjiǎn} — We'll go through security first.
\end{quote}

\begin{cousinweb}
\textbf{\zh{安检}} (\emph{ānjiǎn}) is built from \emph{ān} 'safe' and \zh{检} 'to check', both inherited from Middle Chinese, \emph{'an} and \emph{kjemX} (Guangyun, in Baxter's transcription); it is a short form of a longer name, a safety check.
\end{cousinweb}

\begin{grammarlens}[title={What you've built}]
One more word for getting around.
\end{grammarlens}

\subsection*{Your turn}
\begin{itemize}
\raggedright
  \item \emph{Say it:} \emph{ānjiǎn}
  \item \emph{Say it:} \emph{ānjiǎn}, once more
  \item \emph{Say it:} \emph{wǒmen xiān qù ānjiǎn} — We'll go through security first.
  \item \emph{Recall:} say the Mandarin for to check in a suitcase, then the Mandarin for a boarding pass, then say \emph{ānjiǎn} again
\end{itemize}

\begin{tcolorbox}[breakable,colback=teal!4,colframe=teal!35!black,title={Before you move on}]
What does \zh{安检} mean? (\textquotedblleft{}Security check\textquotedblright{}.) Say it once more.
\end{tcolorbox}
\section[dēngjīkǒu]{dēngjīkǒu — departure gate}

\label{lesson:ZH-C415-dengjikou}



\begin{quote}
Before the new one: say the Mandarin for a boarding pass, then the Mandarin for the security check at an airport.
\end{quote}

\subsection*{You'll want to know: dēngjīkǒu}
\emph{dēngjīkǒu} — \textquotedblleft{}departure gate\textquotedblright{}.

\begin{quote}
\emph{qù shí hào dēngjīkǒu zěnme zǒu} — How do I get to gate ten?
\end{quote}

Its characters come later; for now it is a word for the ear.

\begin{cousinweb}
\emph{dēngjīkǒu} is built from \emph{dēngjī} 'to board a plane' and \zh{口} 'an opening', all inherited from Middle Chinese, \emph{tong}, \emph{kj+j} and \emph{khuwX} (Guangyun, in Baxter's transcription).
\end{cousinweb}

\begin{grammarlens}[title={What you've built}]
One more word for getting around.
\end{grammarlens}

\subsection*{Your turn}
\begin{itemize}
\raggedright
  \item \emph{Say it:} \emph{dēngjīkǒu}
  \item \emph{Say it:} \emph{dēngjīkǒu}, once more
  \item \emph{Say it:} \emph{qù shí hào dēngjīkǒu zěnme zǒu} — How do I get to gate ten?
  \item \emph{Recall:} say the Mandarin for a boarding pass, then the Mandarin for the security check at an airport, then say \emph{dēngjīkǒu} again
\end{itemize}

\begin{tcolorbox}[breakable,colback=teal!4,colframe=teal!35!black,title={Before you move on}]
What does dēngjīkǒu mean? (\textquotedblleft{}Departure gate\textquotedblright{}.) Say it once more.
\end{tcolorbox}
